\begin{abstract}
\addcontentsline{toc}{chapter}{Resumen}
Las organizaciones que ya poseen sensores no logran convertir sus datos en
información utilizable: cada dispositivo emite tramas con formatos distintos, las plataformas
comerciales operan solo con el hardware de su proveedor, y no existe visibilidad del estado de la
conectividad. Este trabajo desarrolla una plataforma web para el registro, monitoreo y gestión de
sensores que recibe tramas de telemetría continua a través de TCP/IP —con MQTT como vía principal y
HTTP como vía secundaria—, las normaliza a un modelo interno único mediante un pipeline de seis
filtros con cuarentena y adaptadores por formato, almacena las mediciones como series de tiempo,
supervisa el estado de cada dispositivo y presenta la información en un panel de control, sobre un caso de validación en el Campus Dr. Luis Rivas del Canto de la Universidad Católica de Temuco. El
contrato nativo de entrada es el estándar abierto SenML (RFC 8428). Esta primera etapa define el
problema mediante un árbol de causas y efectos, formula un objetivo general y cinco específicos con
matriz de trazabilidad, cataloga 68 requerimientos justificados por fuente, evalúa cinco mecanismos
de transporte y cuatro formatos de trama mediante matrices de criterios ponderados —documentando
además el descarte fundamentado del protocolo Contact ID evaluado en la fase previa— y construye un
backlog de 63 elementos cuya verificación numérica delimitó el compromiso al núcleo demostrativo.

{\bf Palabras Clave:}
Internet de las Cosas, monitoreo de sensores, telemetría, MQTT, SenML, normalización de tramas,
series de tiempo, monitoreo de sensores, La Araucanía.
\end{abstract}
\clearpage
